\section{Families below a cofinality}

\import{ordinal.tex}
\import{set/finite.tex}
\import{cardinal/comparison.tex}
\import{cardinal/reconciliation.tex}
\import{cardinal/infinite.tex}
\import{cardinal/regularity.tex}
\import{cardinal/families.tex}
\import{cardinal/cofinality.tex}
\import{set/function-partition.tex}

\begin{theorem}\label{indexed_union_below_cofinality}
    Let $\kappa$ be an infinite cardinal.
    Suppose $A$ is a function on $I$.
    Suppose $\cardLt{I}{\cofinality{\kappa}}$.
    Suppose for every $i\in I$ we have $\cardLt{A(i)}{\kappa}$.
    Then $\cardLt{\unions{\ran{A}}}{\kappa}$.
\end{theorem}
\begin{proof}
    $A$ is a function and $\dom{A}=I$ by definition.
    Hence $\card{\kappa}=\kappa$ by \obv{card_of_cardinal}.
    $\kappa$ is a limit ordinal by
        \obv{infinite_cardinal_is_limit_ordinal}.
    Hence $\kappa$ is an ordinal by definition.
    $\cofinality{\kappa}$ is a cardinal by
        \obv{cofinality_is_cardinal}.
    Hence $\card{\cofinality{\kappa}}=\cofinality{\kappa}$ by
        \obv{card_of_cardinal}.
    $\cofinality{\kappa}\subseteq\kappa$ by
        \obv{cofinality_subseteq_ordinal}.
    \begin{byCase}
        \caseOf{$\kappa=\naturals$.}
            $\naturals$ is a regular cardinal by
                \obv{naturals_are_regular}.
            Hence $\cofinality{\naturals}=\naturals$ by definition.
            Hence $\cofinality{\kappa}=\naturals$.
            Hence $\cardLt{I}{\naturals}$.
            Hence $I$ is finite by \obv{below_omega_iff_finite}.
            $\img{A}{I}$ is finite by
                \obv{finite_function_image,subseteq_refl}.
            $\img{A}{I}=\ran{A}$ by \auto{img_dom}.
            Hence $\ran{A}$ is finite.
            Show that every element of $\ran{A}$ is finite.
            \begin{subproof}
                Fix $B\in\ran{A}$.
                Take $i\in\dom{A}$ such that $A(i)=B$ by
                    \obv{fun_ran_iff}.
                Hence $i\in I$.
                Hence $\cardLt{A(i)}{\kappa}$.
                Hence $\cardLt{B}{\naturals}$.
                Follows by \obv{below_omega_iff_finite}.
            \end{subproof}
            Hence $\unions{\ran{A}}$ is finite by
                \obv{unions_of_finite_family_is_finite}.
            Hence $\cardLt{\unions{\ran{A}}}{\naturals}$ by
                \obv{below_omega_iff_finite}.
        \caseOf{$\kappa\neq\naturals$.}
            Let $s=\{(i,\card{A(i)})\mid i\in I\}$.
            $s$ is a relation by \auto{relation}.
            $s$ is right-unique by \auto{rightunique,pair_eq_iff}.
            Hence $s$ is a function.
            $\dom{s}=I$ by \auto{setext,dom_iff,pair_eq_iff}.
            For every $i\in I$ we have $s(i)=\card{A(i)}$ by
                \obv{function_apply_intro}.
            For every $i\in I$ we have $\card{A(i)}\in\kappa$ by
                \auto{card_elem_iff_cardinal_lt}.
            Hence $s\in\funs{I}{\kappa}$ by \auto{funs_intro}.
            Show that $\lnot\cofinalIn{\ran{s}}{\kappa}$.
            \begin{subproof}
                Suppose $\cofinalIn{\ran{s}}{\kappa}$.
                Hence $\cofinalMap{s}{I}{\kappa}$ by
                    \obv{cofinal_map}.
                Hence $\cofinality{\kappa}\subseteq\card{I}$ by
                    \obv{cofinality_le_domain_card}.
                Hence $\card{\cofinality{\kappa}}\subseteq\card{I}$.
                Hence $\cardLe{\cofinality{\kappa}}{I}$ by
                    \obv{card_subseteq_iff_cardinal_le}.
                $\cardLe{I}{\cofinality{\kappa}}$ and
                    $I\not\approx\cofinality{\kappa}$ by
                    \obv{cardinal_lt}.
                Hence $I\approx\cofinality{\kappa}$ by
                    \obv{cardinal_le_antisymmetric_up_to_equinum}.
                Contradiction.
            \end{subproof}
            $\ran{s}\subseteq\kappa$ by \obv{funs_ran}.
            Every element of $\ran{s}$ is an ordinal by
                \obv{subseteq,elem_of_ordinal_is_ordinal,limit_ordinal}.
            Hence $\unions{\ran{s}}$ is an ordinal by
                \obv{union_of_ordinals_is_ordinal}.
            Let $u=\unions{\ran{s}}$.
            $u$ is an ordinal.
            Show that $u\subseteq\kappa$.
            \begin{subproof}
                Show that every element of $u$ is an element of $\kappa$.
                \begin{subproof}
                    Fix $x\in u$.
                    Take $\sigma\in\ran{s}$ such that $x\in\sigma$ by
                        \obv{unions_iff}.
                    Hence $\sigma\in\kappa$ by \obv{subseteq}.
                    Follows by \obv{ordinal_transitivity,limit_ordinal}.
                \end{subproof}
                Follows by \obv{subseteq}.
            \end{subproof}
            $u\neq\kappa$ by
                \obv{cofinal_in_limit_iff_union}.
            Hence $u\in\kappa$ by
                \obv{ordinal_subseteq_implies_elem_or_eq}.
            $\suc{u}\in\kappa$ by \obv{suc_elem_limitordinal}.
            $u\subseteq\suc{u}$ by \obv{subseteq_self_suc_intro}.
            Let $\delta=\suc{u}$.
            $\delta\in\kappa$ and $u\subseteq\delta$.
            $\card{I}\in\card{\cofinality{\kappa}}$ by
                \obv{card_elem_iff_cardinal_lt}.
            Hence $\card{I}\in\cofinality{\kappa}$.
            Hence $\card{I}\in\kappa$ by \obv{elem_subseteq}.
            $\naturals\subseteq\kappa$ by
                \obv{naturals_subseteq_infinite_cardinal}.
            $\naturals$ is an ordinal by \obv{omega_is_ordinal}.
            Hence $\naturals\in\kappa$ by
                \obv{ordinal_subseteq_implies_elem_or_eq}.
            $(\delta\union\card{I})\union\naturals\in\kappa$ by
                \obv{three_ordinal_union_below}.
            $\naturals\subseteq
                (\delta\union\card{I})\union\naturals$ by
                \obv{union_upper_right}.
            $\card{I}\subseteq
                (\delta\union\card{I})\union\naturals$ by
                \auto{union_upper_left,union_upper_right,
                    subseteq_transitive}.
            $\delta\subseteq
                (\delta\union\card{I})\union\naturals$ by
                \auto{union_upper_left,subseteq_transitive}.
            Let $\gamma=(\delta\union\card{I})\union\naturals$.
            $\gamma\in\kappa$.
            $\naturals\subseteq\gamma$ and $\card{I}\subseteq\gamma$ and
                $\delta\subseteq\gamma$.
            $\card{\gamma}\in\kappa$ by
                \obv{card_of_ordinal_below}.
            $\card{\gamma}$ is a cardinal by \obv{card_is_cardinal}.
            $\gamma\approx\card{\gamma}$ by
                \obv{equinumerous_card}.
            Hence $\cardLe{\gamma}{\card{\gamma}}$ by
                \auto{equinum_cardinal_le}.
            Let $\theta=\card{\gamma}$.
            $\theta\in\kappa$ and $\theta$ is a cardinal.
            $\cardLe{\gamma}{\theta}$.
            Hence $\cardLe{\naturals}{\gamma}$ by
                \obv{subseteq_cardinal_le}.
            Hence $\card{\naturals}\subseteq\card{\gamma}$ by
                \auto{card_subseteq_iff_cardinal_le}.
            $\card{\naturals}=\naturals$ by \obv{card_of_naturals}.
            Hence $\naturals\subseteq\theta$.
            $\naturals$ is infinite by \obv{naturals_is_infinite}.
            Hence $\theta$ is infinite by \obv{infinite_superset}.
            Hence $\cardLe{\card{I}}{\gamma}$ by
                \obv{subseteq_cardinal_le}.
            $I\approx\card{I}$ by \obv{equinumerous_card}.
            Hence $\cardLe{I}{\card{I}}$ by
                \obv{equinum_cardinal_le}.
            Hence $\cardLe{I}{\gamma}$ by
                \obv{cardinal_le_transitive}.
            Hence $\cardLe{I}{\theta}$ by
                \obv{cardinal_le_transitive}.
            Show that for every $i\in I$ we have
                $\cardLe{A(i)}{\theta}$.
            \begin{subproof}
                Fix $i\in I$.
                $s(i)=\card{A(i)}$.
                Hence $\card{A(i)}\in\ran{s}$ by
                    \auto{fun_ran_iff,funs}.
                Hence $\card{A(i)}\subseteq u$ by
                    \auto{subseteq,unions_iff}.
                Hence $\card{A(i)}\subseteq\delta$ by
                    \obv{subseteq_transitive}.
                Show that $\card{A(i)}\subseteq\gamma$.
                \begin{subproof}
                    Show that every element of $\card{A(i)}$ is an element
                        of $\gamma$.
                    \begin{subproof}
                        Fix $x\in\card{A(i)}$.
                        Hence $x\in\delta$ by \obv{subseteq}.
                        Hence $x\in\delta\union\card{I}$ by
                            \obv{union_intro_left}.
                        Hence $x\in(\delta\union\card{I})\union\naturals$ by
                            \obv{union_intro_left}.
                    \end{subproof}
                    Follows by \obv{subseteq}.
                \end{subproof}
                Hence $\cardLe{\card{A(i)}}{\gamma}$ by
                    \obv{subseteq_cardinal_le}.
                $A(i)\approx\card{A(i)}$ by
                    \obv{equinumerous_card}.
                Hence $\cardLe{A(i)}{\card{A(i)}}$ by
                    \obv{equinum_cardinal_le}.
                Follows by \obv{cardinal_le_transitive}.
            \end{subproof}
            Hence $\cardLe{\unions{\ran{A}}}{\theta}$ by
                \obv{indexed_union_infinite_cardinal_bound}.
            $\cardLt{\theta}{\kappa}$ by
                \obv{cardinal_element_small}.
            Follows by \obv{cardinal_le_lt_transitive}.
    \end{byCase}
\end{proof}

\begin{corollary}\label{regular_indexed_union}
    Let $\kappa$ be a regular cardinal.
    Suppose $A$ is a function on $I$.
    Suppose $\cardLt{I}{\kappa}$.
    Suppose for every $i\in I$ we have $\cardLt{A(i)}{\kappa}$.
    Then $\cardLt{\unions{\ran{A}}}{\kappa}$.
\end{corollary}
\begin{proof}
    $\kappa$ is an infinite cardinal and
        $\cofinality{\kappa}=\kappa$ by \obv{regular_cardinal}.
    Hence $\cardLt{I}{\cofinality{\kappa}}$.
    Follows by \obv{indexed_union_below_cofinality}.
\end{proof}

\begin{theorem}\label{small_binary_union}
    Let $\kappa$ be an infinite cardinal.
    Suppose $\cardLt{A}{\kappa}$ and $\cardLt{B}{\kappa}$.
    Then $\cardLt{A\union B}{\kappa}$.
\end{theorem}
\begin{proof}
    $\kappa$ is a limit ordinal by
        \obv{infinite_cardinal_is_limit_ordinal}.
    Hence $\cofinality{\kappa}$ is an infinite cardinal by
        \obv{cofinality_of_limit_is_infinite_cardinal}.
    $\{A,B\}$ is finite by \obv{pair_is_finite}.
    Hence $\cardLt{\{A,B\}}{\cofinality{\kappa}}$ by
        \obv{finite_below_infinite}.
    $\identity{\{A,B\}}$ is a function on $\{A,B\}$ by
        \obv{id_is_function_on}.
    For every $C\in\{A,B\}$ we have $\cardLt{C}{\kappa}$ by
        \auto{pairset_iff}.
    Show that for every $C\in\{A,B\}$ we have
        $\cardLt{(\identity{\{A,B\}})(C)}{\kappa}$.
    \begin{subproof}
        Fix $C\in\{A,B\}$.
        $(\identity{\{A,B\}})(C)=C$ by \obv{id_apply}.
    \end{subproof}
    Hence $\cardLt{\unions{\ran{\identity{\{A,B\}}}}}{\kappa}$ by
        \obv{indexed_union_below_cofinality}.
    $\ran{\identity{\{A,B\}}}=\{A,B\}$ by \obv{id_ran}.
    $\unions{\{A,B\}}=A\union B$ by \obv{union_as_unions}.
\end{proof}

\begin{corollary}\label{small_cons}
    Let $\kappa$ be an infinite cardinal.
    Suppose $\cardLt{A}{\kappa}$.
    Then $\cardLt{\cons{x}{A}}{\kappa}$.
\end{corollary}
\begin{proof}
    $\{x\}$ is finite by \obv{singleton_is_finite}.
    Hence $\cardLt{\{x\}}{\kappa}$ by \obv{finite_below_infinite}.
    Hence $\cardLt{\{x\}\union A}{\kappa}$ by
        \obv{small_binary_union}.
    $\{x\}\union A=\cons{x}{A}$ by
        \auto{setext,singleton_iff,union_iff,cons_iff}.
\end{proof}

\begin{theorem}\label{large_minus_small}
    Let $\kappa$ be an infinite cardinal.
    Suppose $X\approx\kappa$ and $\cardLt{F}{\kappa}$.
    Then $X\setminus F\approx\kappa$.
\end{theorem}
\begin{proof}
    $X\setminus F\subseteq X$ by \obv{setminus_subseteq}.
    Hence $\cardLe{X\setminus F}{X}$ by
        \obv{subseteq_cardinal_le}.
    $\cardLe{X}{\kappa}$ by \auto{equinum_cardinal_le}.
    Hence $\cardLe{X\setminus F}{\kappa}$ by
        \obv{cardinal_le_transitive}.
    Show that $X\setminus F\approx\kappa$.
    \begin{subproof}
        Suppose not.
        Hence $\cardLt{X\setminus F}{\kappa}$ by
            \obv{cardinal_lt}.
        $X\inter F\subseteq F$ by \obv{inter_lower_right}.
        Hence $\cardLe{X\inter F}{F}$ by
            \obv{subseteq_cardinal_le}.
        Hence $\cardLt{X\inter F}{\kappa}$ by
            \obv{cardinal_le_lt_transitive}.
        Hence $\cardLt{(X\setminus F)\union(X\inter F)}{\kappa}$ by
            \obv{small_binary_union}.
        $X=(X\setminus F)\union(X\inter F)$ by
            \auto{setext,setminus,inter,union_iff}.
        Hence $\cardLt{X}{\kappa}$.
        Hence $\lnot(X\approx\kappa)$ by \obv{cardinal_lt}.
        Contradiction.
    \end{subproof}
\end{proof}

\begin{theorem}\label{large_fibre_below_cofinality}
    Let $\kappa$ be an infinite cardinal.
    Suppose $K\approx\kappa$.
    Let $f\in\funs{K}{J}$.
    Suppose $\cardLt{J}{\cofinality{\kappa}}$.
    Then there exists $j\in J$ such that
        $\fibre{f}{j}\approx\kappa$.
\end{theorem}
\begin{proof}
    $f$ is a function and $\dom{f}=K$ by
        \obv{funs_is_function,funs}.
    $\indexedFibres{f}{J}\in\funs{J}{\pow{\dom{f}}}$ by
        \obv{indexed_fibres_type}.
    Hence $\indexedFibres{f}{J}$ is a function and
        $\dom{\indexedFibres{f}{J}}=J$ by
        \obv{funs_is_function,funs}.
    Hence $\indexedFibres{f}{J}$ is a function on $J$.
    Suppose not.
    Hence for every $j\in J$ we have
        $\lnot(\fibre{f}{j}\approx\kappa)$ by \auto.
    Show that for every $j\in J$ we have
        $\cardLt{\fibre{f}{j}}{\kappa}$.
    \begin{subproof}
        Fix $j\in J$.
        $\fibre{f}{j}\subseteq\dom{f}$ by
            \obv{subseteq,function_fibre_iff}.
        Hence $\fibre{f}{j}\subseteq K$.
        Hence $\cardLe{\fibre{f}{j}}{K}$ by
            \obv{subseteq_cardinal_le}.
        $\cardLe{K}{\kappa}$ by \auto{equinum_cardinal_le}.
        Hence $\cardLe{\fibre{f}{j}}{\kappa}$ by
            \obv{cardinal_le_transitive}.
        Follows by \obv{cardinal_lt}.
    \end{subproof}
    Show that for every $j\in J$ we have
        $\cardLt{(\indexedFibres{f}{J})(j)}{\kappa}$.
    \begin{subproof}
        Fix $j\in J$.
        $(\indexedFibres{f}{J})(j)=\fibre{f}{j}$ by
            \obv{indexed_fibres_apply}.
        Follows by \auto.
    \end{subproof}
    Hence $\cardLt{\unions{\ran{\indexedFibres{f}{J}}}}{\kappa}$ by
        \obv{indexed_union_below_cofinality}.
    $\unions{\ran{\indexedFibres{f}{J}}}=K$ by
        \obv{indexed_fibres_cover_domain}.
    Hence $\cardLt{K}{\kappa}$.
    Hence $\lnot(K\approx\kappa)$ by \obv{cardinal_lt}.
    Contradiction.
\end{proof}

\begin{corollary}\label{regular_large_fibre}
    Let $\rho$ be a regular cardinal.
    Suppose $K\approx\rho$.
    Let $f\in\funs{K}{J}$.
    Suppose $\cardLt{J}{\rho}$.
    Then there exists $j\in J$ such that
        $\fibre{f}{j}\approx\rho$.
\end{corollary}
\begin{proof}
    $\rho$ is an infinite cardinal and
        $\cofinality{\rho}=\rho$ by \obv{regular_cardinal}.
    Hence $\cardLt{J}{\cofinality{\rho}}$.
    Follows by \obv{large_fibre_below_cofinality}.
\end{proof}
